
\subsection{Theoretical Results}
\label{sec:theory}


We now present our main theoretical contributions.
In order to simplify the exposition and convey the gist of our results, we will
only present a simplified version of our theorems.
We will suppress constants, $\polylog$ terms, and other technical
details that arise due to a covering argument in our proofs.
A rigorous treatment is available in Appendix~\ref{sec:analysis}.


\textbf{Maximum Information Gain:}
Up until this point, we have not discussed much about the kernel $\phix$ of the
domain $\Xcal$.
Since we are optimising $\func$ over $\Xcal$, it is natural to expect that this
will appear in the bounds.
\citet{srinivas10gpbandits} showed that the statistical difficulty of GP bandits is
determined by the \emph{Maximum Information Gain} (\mig) which measures the maximum
information a subset of observations have about $\func$. 
We denote it by $\IGn(A)$ where $A$ is a subset of $\Xcal$ and $n$ is the number of
queries to $\func$. We refer the reader to Appendix~\ref{sec:analysis} for a formal definition of \migs. For the current exposition however, it suffices to know that 
$\IGn(A)$ depends
on the domain kernel $\phix$, the number of times $n$ we have queried $\func$, and the
volume $\vol(A)$ of the set $A\subset\Xcal$. 
The latter dependence on $\vol(A)$ will be most important to us.
For instance, when we use an SE kernel for $\phix$, we have
$\IGn(A) \propto \vol(A)\,\logn^{d+1}$%
~\citep{seeger08information}.
~\citet{srinivas10gpbandits} showed that 
the simple regret $S(\capital)$ for \gpucbs after capital $\capital$ can be bounded by,
\begin{align*}
\hspace{-0.1in}
\text{Simple Regret for \gpucb:}\hspace{0.15in}
S(\capital) \lesssim 
  \sqrt{\frac{\IGnc(\Xcal)}{\ncapital}},
\label{eqn:gpucbregret}
\numberthis
\end{align*}
where $\ncapital = \floor{\capital/\cost(\zhf)}$.


In our analysis of \mfgpucbtwos we show that most queries to $\gunc$ at fidelity
$\zhf$ will be confined to a small subset of the domain $\Xcal$ which contains the
 optimum $\xopt$.
More precisely, after capital $\capital$, for any $\alpha\in(0,1)$, we show that 
there exists $\rho>0$ such that the number of queries \emph{outside} the following
set $\Xcalrho$ is less than $\ncapital^\alpha$. 
\begin{align*}
\Xcalrho = \big\{x\in\Xcal\,:\, \funcopt - \func(x)
\leq 2\rho\sqrt{\kernelscale}\,\inffunc(\diamz) \big\}.
\label{eqn:Xcalrho}
\numberthis
\end{align*}
Here, $\inffunc$ is from~\eqref{eqn:inffunc} and $\diamz$ is the $L_2$
diameter of $\Zcal$. %\todo{the following lines are not very clear.}
While it is true that any optimisation algorithm would eventually query extensively
in a neighbourhood around the optimum, a strong result of the above form is
not always possible.
For instance, in the case of \gpucbs, the best achievable bound on the number of queries
in any set that does not contain $\xopt$ is $\ncapital^{1/2}$. 
The fact that the above set $\Xcalrho$ exists relies crucially on the multi-fidelity
assumptions and the fact that our algorithm leverages information from lower
fidelities when querying at $\zhf$.
% and is determined by the smoothness of $\gunc$ across $\Zcal$.
As $\inffunc$ is small when $\gunc$ is smooth across $\Zcal$,  the set $\Xcalrho$ will
be small when the approximations are highly informative about $\gunc(\zhf, \cdot)$. 
To see this more clearly, consider again the case where $\phiz$ is a SE kernel, 
where we have 
 $\Xcalrho \approx \{x\in\Xcal: \funcopt - \func(x) \leq 2\rho\sqrt{p}/\hz\}$.
When $\hz$ is large and $\gunc$ is smooth across $\Zcal$, $\Xcalrho$ is small 
as the right side of the inequality is smaller.
As \mfgpucbtwos confines most of its evaluations to this small set
containing $\xopt$, we will be able to achieve much better regret than \gpucb.
When $\hz$ is small and $\gunc$ is not smooth across $\Zcal$,
the set $\Xcalrho$ becomes large and the advantage of multi-fidelity
optimisation diminishes.


We now provide an informal statement of our main result below.
$\lesssim,\asymp$ will denote inequality and equality ignoring constant and
$\polylog$ terms.

\begin{theorem}[Informal, Regret of \mfgpucbtwo]
\label{thm:regret}
Let $g\sim\GP(0,\kernel)$ where $\kernel$ satisfies~\eqref{eqn:mfkernel}.
Choose $\betat \asymp d\log(t/\delta)$.
Then, for sufficiently large $\capital$ and for all $\alpha \in(0,1)$, there exists $\rho$
depending on $\alpha$ such that the following bound holds w.h.p.
\[
S(\capital) \;\lesssim\;
  \sqrt{\frac{\IGnc(\Xcalrho)}{\ncapital}}
  \;+\;
  \sqrt{\frac{\IGncalpha(\Xcal)}{\ncapital^{2-\alpha}}}
\]
\end{theorem}
In the above bound, the latter term vanishes fast due to the $\ncapital^{-(1-\alpha/2)}$
dependence.
When comparing this with~\eqref{eqn:gpucbregret},
we see that we outperform \gpucbs by a factor of
$\sqrt{\IGnc(\Xcalrho)/\IGnc(\Xcal)} \asymp \sqrt{\vol(\Xcalrho)/\vol(\Xcal)}$
asymptotically.
If $\gunc$ is smooth across the fidelity space, $\Xcalrho$ is small and
the gains over \gpucbs are significant.
If $\gunc$ becomes less smooth across $\Zcal$, the bound decays gracefully, but
we are never worse than \gpucbs up to constant factors.

Theorem~\ref{thm:regret} also has similarities to the bounds 
of~\citet{kandasamy16mfbo} who also demonstrate better regret than \gpucbs by
showing that it is dominated by queries inside a set $\Xcal'$ which contains the optimum.
However, their bounds  depend critically on certain threshold hyper-parameters
which determine the volume of $\Xcal'$ among other  terms in their regret.
The authors of that paper note that their bounds will suffer if these hyper-parameters are not
chosen appropriately, but do not provide theoretically justified methods to make this choice.
In contrast, many of the design choices for \bocas fall out naturally of 
our modeling assumptions.
Beyond this analogue, our results are not comparable to~\citet{kandasamy16mfbo}
as the assumptions are different.

